\documentclass[10pt,a4paper]{amsart}
\usepackage{amsmath,amssymb,amsthm,graphicx}
% Retained minimal preamble for the Stacks Project chapter more-morphisms.tex (More on Morphisms).
% Verbatim excerpt of the chapter preamble: Xy-pic setup (native preamble lines 14-18) and the native
% theorem declarations with the mathematical macros (native preamble lines 154-217). Copied mechanically;
% nothing added, removed or redefined. The wrapper class and the amsmath/amssymb/amsthm/graphicx base
% packages are supplied by the assembler, so the class/externaldocument part of the native preamble is omitted.
% For commutative diagrams we use Xy-pic
\usepackage[all]{xy}

% We use 2cell for 2-commutative diagrams.
\xyoption{2cell}
% Theorem environments.
%
\theoremstyle{plain}
\newtheorem{theorem}[subsection]{Theorem}
\newtheorem{proposition}[subsection]{Proposition}
\newtheorem{lemma}[subsection]{Lemma}

\theoremstyle{definition}
\newtheorem{definition}[subsection]{Definition}
\newtheorem{example}[subsection]{Example}
\newtheorem{exercise}[subsection]{Exercise}
\newtheorem{situation}[subsection]{Situation}

\theoremstyle{remark}
\newtheorem{remark}[subsection]{Remark}
\newtheorem{remarks}[subsection]{Remarks}

\numberwithin{equation}{subsection}

% Macros
%
\def\lim{\mathop{\mathrm{lim}}\nolimits}
\def\colim{\mathop{\mathrm{colim}}\nolimits}
\def\Spec{\mathop{\mathrm{Spec}}}
\def\Hom{\mathop{\mathrm{Hom}}\nolimits}
\def\Ext{\mathop{\mathrm{Ext}}\nolimits}
\def\SheafHom{\mathop{\mathcal{H}\!\mathit{om}}\nolimits}
\def\SheafExt{\mathop{\mathcal{E}\!\mathit{xt}}\nolimits}
\def\Sch{\mathit{Sch}}
\def\Mor{\mathop{\mathrm{Mor}}\nolimits}
\def\Ob{\mathop{\mathrm{Ob}}\nolimits}
\def\Sh{\mathop{\mathit{Sh}}\nolimits}
\def\NL{\mathop{N\!L}\nolimits}
\def\CH{\mathop{\mathrm{CH}}\nolimits}
\def\proetale{{pro\text{-}\acute{e}tale}}
\def\etale{{\acute{e}tale}}
\def\QCoh{\mathit{QCoh}}
\def\Ker{\mathop{\mathrm{Ker}}}
\def\Im{\mathop{\mathrm{Im}}}
\def\Coker{\mathop{\mathrm{Coker}}}
\def\Coim{\mathop{\mathrm{Coim}}}

% Boxtimes
%
\DeclareMathSymbol{\boxtimes}{\mathbin}{AMSa}{"02}

%
% Macros for moduli stacks/spaces
%
\def\QCohstack{\mathcal{QC}\!\mathit{oh}}
\def\Cohstack{\mathcal{C}\!\mathit{oh}}
\def\Spacesstack{\mathcal{S}\!\mathit{paces}}
\def\Quotfunctor{\mathrm{Quot}}
\def\Hilbfunctor{\mathrm{Hilb}}
\def\Curvesstack{\mathcal{C}\!\mathit{urves}}
\def\Polarizedstack{\mathcal{P}\!\mathit{olarized}}
\def\Complexesstack{\mathcal{C}\!\mathit{omplexes}}
% \Pic is the operator that assigns to X its picard group, usage \Pic(X)
% \Picardstack_{X/B} denotes the Picard stack of X over B
% \Picardfunctor_{X/B} denotes the Picard functor of X over B
\def\Pic{\mathop{\mathrm{Pic}}\nolimits}
\def\Picardstack{\mathcal{P}\!\mathit{ic}}
\def\Picardfunctor{\mathrm{Pic}}
\def\Deformationcategory{\mathcal{D}\!\mathit{ef}}
\newenvironment{reference}{}{}
\newenvironment{slogan}{}{}
\newenvironment{history}{}{}
\usepackage{cleveref}
\crefname{originallemma}{original lemma}{original lemmas}
\Crefname{originallemma}{Original lemma}{Original lemmas}
\makeatletter
\newlabel{lemma-quasi-finite-quasi-section-meeting-nearby-open}{{05GT}{source}{}{sourcepointer-1}{}}
\newlabel{lemma-quasi-finite-quasi-section-meeting-nearby-open@cref}{{[originallemma][0][]05GT}{[1][0][]source}{}{}{}}
\makeatother
\title{Global form of the nearby-open quasi-section: a curve through x inside U quasi-finite over the base at x}
\author{The Stacks Project Authors}
\date{stacks-project:more-morphisms.tex:local-snapshot-20261011; GNU FDL 1.2 or later}
\begin{document}
\maketitle
\noindent\textit{Editorial note.} Statement span is lines 14993-15004 of sources/stacks-more-morphisms/pinned-original-source.tex; proof span is lines 15006-15023, the complete author proof (16 body lines).\par\medskip
\noindent\textit{Editorial note.} Navigation stubs: 1 cited-but-not-included labels are supplied as non-load-bearing pointers with their pinned tags where the snapshot tag index has them.\par\medskip
\noindent\textit{Editorial note.} extra\_packages is empty on purpose: the retained preamble already loads xy with options and a second declaration clashes in the assembler wrapper.\par\medskip
\noindent\textit{Editorial note.} Text is retained verbatim, including the author's own typos and placeholders.\par\medskip
\noindent\textit{Editorial note.} This is the X-version cited-but-not-yet-opened when L16095 (change-hypotheses) was collected last wave; the citation chain in the lane is now closed in the right order. Sibling of L16081 (meeting-nearby-open without the X suffix): different hypothesis packaging, both counted per their distinct aggregate statements.\par\medskip
\noindent\textit{Allowed standard prerequisites (editorial).}\begin{itemize}
\item Explicit prerequisite boundary. The chapter's local lemma quasi-finite-quasi-section-meeting-nearby-open (the affine-local algebraic version collected as an earlier wave card); elementary properties of local rings of schemes at a point and of reduced closures.
\end{itemize}\medskip
\ifdefined\hypertarget\hypertarget{sourcepointer-1}{}\fi
\noindent\textit{Source reference.} Original lemma 05GT: sources/stacks-more-morphisms/pinned-original-source.tex, source lines 14958--14967. Cited by the retained statement or proof but not part of the collected spans; supplied only as a navigation pointer. The argument does not depend on its content.\par\medskip
\begin{lemma}
\label{lemma-quasi-finite-quasi-section-meeting-nearby-open-X}
Let $f : X \to S$ be a morphism of schemes.
Let $x \in X$ with image $s \in S$.
Let $U \subset X$ be an open subscheme.
Assume $f$ locally of finite type, $S$ locally Noetherian, $x$ a closed
point of $X_s$, and assume there exists a point $x' \in U$ with
$x' \leadsto x$ and $f(x') \not = s$. Then there exists a closed
subscheme $Z \subset X$ such that (a) $x \in Z$, (b) $f|_Z : Z \to S$ is
quasi-finite at $x$, and (c) there exists a $z \in Z$, $z \in U$,
$z \leadsto x$ and $f(z) \not = s$.
\end{lemma}

\begin{proof}
This is a reformulation of
Lemma \ref{lemma-quasi-finite-quasi-section-meeting-nearby-open}.
Namely, set $A = \mathcal{O}_{S, s}$ and $B = \mathcal{O}_{X, x}$.
Denote $V \subset \Spec(B)$ the inverse image of $U$.
The ring map $f^\sharp : A \to B$ is essentially of finite type.
By assumption there exists at least one point of $V$ which does not
map to the closed point of $\Spec(A)$. Hence all the assumptions of
Lemma \ref{lemma-quasi-finite-quasi-section-meeting-nearby-open}
hold and we obtain a prime $\mathfrak q \subset B$ which does not
lie over $\mathfrak m_A$ and such that $A \to B/\mathfrak q$ is
the localization of a quasi-finite ring map. Let $z \in X$ be the
image of the point $\mathfrak q$ under the canonical
morphism $\Spec(B) \to X$. Set $Z = \overline{\{z\}}$
with the induced reduced scheme structure. As $z \leadsto x$
we see that $x \in Z$ and $\mathcal{O}_{Z, x} = B/\mathfrak q$.
By construction $Z \to S$ is quasi-finite at $x$.
\end{proof}

\end{document}
